\documentclass[12pt]{article}
\usepackage[utf8]{inputenc}
\usepackage{geometry}
\geometry{a4paper, margin=1in}
\usepackage{graphicx}
\usepackage{amsmath}
\usepackage{hyperref}

\title{Multimodal 3D Gaussian Splatting:\\Integrating Near-Infrared Features via Vision Transformers (ViT)}
\author{}
\date{\today}

\begin{document}

\maketitle

\begin{abstract}
This report details the implementation of a Multimodal 3D Gaussian Splatting (3DGS) pipeline designed to fuse high-resolution RGB data with Near-Infrared (NIR) data. To avoid the geometric distortions introduced by traditional 2D image registration, we embed the NIR data directly into the 3D space. Building upon recent advancements in representation learning, we employ a Vision Transformer (ViT) as a frozen feature encoder to extract rich, semantic targets from the NIR images. We detail the modified 3DGS architecture, the distillation loss functions, and discuss current limitations stemming from data collection methodologies.
\end{abstract}

\section{Introduction}
Combining multiple modalities, such as visible spectrum (RGB) and Near-Infrared (NIR), provides a wealth of information for scene understanding. Traditionally, combining a 4K RGB image with a 1080p NIR image requires 2D image registration (warping and stretching). This process is highly error-prone and severely degrades the high-frequency geometric details of the RGB data.

To resolve this, we propose a ``True 3D Fusion'' approach. By modifying the core data structure of 3D Gaussian Splatting, we allow the 3D geometry to hold NIR information intrinsically. Consequently, the modalities are perfectly aligned in 3D space without requiring 2D pixel warping.

\section{Methodology}

\subsection{Joint 3D Alignment (Structure-from-Motion)}
The pipeline begins by ingesting both raw RGB and raw NIR images into a joint Structure-from-Motion (SfM) pipeline using COLMAP. The algorithm extracts SIFT keypoints across all images indiscriminately. By matching spatial features between the modalities, COLMAP estimates a unified set of 3D camera poses and generates a shared Sparse Point Cloud. This point cloud provides the initial spatial coordinates (X, Y, Z) for the 3D Gaussians.

\subsection{Multimodal 3D Gaussians}
Standard 3D Gaussians are parameterized by position, scale, rotation, opacity, and spherical harmonics (representing RGB color). To support NIR data, we expand the Gaussian data structure. Every Gaussian is augmented with an additional 32-dimensional feature vector, denoted as $f_{nir}$. This high-dimensional vector gives the model the mathematical capacity (or ``vocabulary'') to represent complex spatial textures, avoiding the representation collapse that often occurs when learning single-channel grayscale values.

\subsection{Vision Transformer (ViT) Target Generation}
To guide the Gaussians in learning the NIR data, we utilize a Vision Transformer (ViT). ViTs (such as DINO or standard ViT architectures) have shown exceptional capability in extracting deep structural and semantic representations from images, heavily outperforming traditional CNNs in capturing global context.

During the training loop, when an NIR camera pose is selected, the corresponding raw NIR image is passed through the frozen ViT. Since ViTs often output very high-dimensional embeddings (e.g., 384 or 768 dimensions) which would cause GPU memory exhaustion if stored per-Gaussian, a lightweight linear projection head maps the ViT output down to our target 32 dimensions. This 32D feature map serves as the ``ground truth'' structural target that the Gaussians must learn to replicate.

\subsection{Rendering Architecture}
Because the original Inria CUDA rasterizer is hardcoded to render exactly three color channels, we transitioned the rendering backend to \texttt{gsplat}. The \texttt{gsplat} engine natively supports the forward-projection and rasterization of arbitrary high-dimensional arrays, allowing us to render 2D maps of our 32D $f_{nir}$ vectors.

\subsection{Masked Feature Distillation Loss}
When training on an NIR view, the system compares the rendered 32D $f_{nir}$ map against the 32D target map generated by the ViT. 
Because we are comparing high-dimensional vectors rather than absolute pixel intensities, we employ \textbf{Cosine Similarity} as the loss function. 

Furthermore, to prevent the Gaussians from hallucinating NIR features in empty space (e.g., the sky), we explicitly mask the loss using the rendered opacity map ($\alpha$-map):
\begin{equation}
\mathcal{L}_{NIR} = \alpha\_map \times (1 - \text{CosineSimilarity}(Rendered\_f_{nir}, Target\_ViT_{nir}))
\end{equation}
This ensures that the model is only penalized for errors on solid geometry. For RGB views, standard $L_1$ and D-SSIM loss functions are used against the real RGB photos.

\section{Implementation Details}
A critical challenge when transitioning to the \texttt{gsplat} backend was the disruption of the gradient flow required for Gaussian densification. 3DGS relies on gradients to identify blurry areas and clone or split Gaussians accordingly. We engineered a solution by manually extracting the absolute gradient (\texttt{.absgrad}) from the \texttt{gsplat} backend during the backward pass and routing it into the adaptive density control logic, successfully restoring the model's ability to grow.

\section{Discussion and Limitations}
While the software architecture and ViT-based distillation successfully compiled and minimized the masked loss during training (RGB $L_1$ dropped to 0.0414; NIR Cosine dropped to 0.0118), visual validation highlighted a critical flaw in the data collection methodology.

Because the RGB and NIR videos were captured sequentially and with significantly different fields of view (zoom levels), the COLMAP SfM algorithm failed to align the NIR cameras correctly with the RGB cameras. Consequently, COLMAP placed the NIR cameras pointing into empty 3D space. 

Our masking logic functioned perfectly: because the cameras viewed empty space, the $\alpha$-map was 0. This artificially drove the loss to near-zero and resulted in flat, unlearned feature vectors (visualized as a ''green blob'' in PCA). To achieve visually accurate results in future iterations, the dataset must be captured using a synchronized, stereo-calibrated rig where RGB and NIR cameras trigger simultaneously with similar fields of view.

\section{Conclusion}
We successfully designed and implemented a Multimodal 3D Gaussian Splatting architecture. By utilizing a Vision Transformer for target generation and projecting the high-dimensional features into the 3D Gaussians, we achieved true 3D fusion. Once the data collection methodology is updated to provide synchronized poses, this pipeline will yield highly accurate, multi-spectral 3D reconstructions.

\end{document}
